% Generated draft: requires mathematical and source review.
\documentclass[11pt]{article}
\usepackage[margin=25mm]{geometry}
\usepackage{fontspec}
\setmainfont{Latin Modern Roman}
\usepackage{amsmath,amssymb,mathtools,unicode-math}
\setmathfont{Latin Modern Math}
\usepackage{enumitem,needspace,xurl,hyperref}
\setlength{\parindent}{0pt}
\setlength{\parskip}{5pt}
\setlength{\emergencystretch}{3em}
\newcommand{\N}{\mathbb{N}}
\newcommand{\Z}{\mathbb{Z}}
\newcommand{\Q}{\mathbb{Q}}
\newcommand{\R}{\mathbb{R}}
\newcommand{\C}{\mathbb{C}}
\newcommand{\F}{\mathbb{F}}
\newcommand{\A}{\mathbb{A}}
\newcommand{\PP}{\mathbb{P}}
\newcommand{\E}{\mathbb{E}}
\newcommand{\eps}{\varepsilon}
\newcommand{\dd}{\,\mathrm{d}}
\begin{document}
\section*{UnsolvedMath ID: 30004330}
\textbf{Source problem code:} catalogue entry
\subsection{Definitions and mathematical statement}
The supplied record does not specify a ground field or its characteristic. The following notation gives a conditional formulation over a fixed algebraically closed field \(k\); it does not assert that every characteristic is intended.

Let \(C\subset\PP^2_k\) be a reduced plane curve, and let \(Z=\operatorname{Sing}(C)_{\mathrm{red}}\) be its finite set of singular points, endowed with the reduced structure. For \(Z\ne\varnothing\), put \(S=k[x_0,x_1,x_2]\) and let
\[
I_Z=\bigcap_{p\in Z}\mathfrak p_p\subset S,
\qquad
I_Z^{(m)}=\bigcap_{p\in Z}\mathfrak p_p^{\,m}
\quad(m\in\Z_{\ge1}),
\]
where \(\mathfrak p_p\) is the homogeneous ideal of the point \(p\). Thus \(I_Z^{(m)}\) consists of homogeneous forms vanishing to order at least \(m\) at every point of \(Z\). For a nonzero homogeneous ideal \(J\subset S\), define its initial degree by
\[
\alpha(J)=\min\{d\in\Z_{\ge0}:J_d\ne0\}.
\]
The Waldschmidt constant is
\[
\widehat\alpha(I_Z)
=\lim_{m\to\infty}\frac{\alpha(I_Z^{(m)})}{m}
=\inf_{m\ge1}\frac{\alpha(I_Z^{(m)})}{m}.
\]

The source asks whether, over its intended ground field, every reduced plane curve with nonempty singular set satisfies
\[
\widehat\alpha\bigl(I_{\operatorname{Sing}(C)_{\mathrm{red}}}\bigr)\in\Q,
\]
and whether this constant can be computed. It specifically asks for its computation when \(C\) is the union of the \(21\) lines in Klein's arrangement, with \(Z\) the full singular set of that union.

The ground field, the meaning of computability, and a defining realization of the named arrangement are absent from the supplied record. Consequently, the computational request is retained as a request to determine the constant, without imposing an unspecified algorithmic model. The source also does not address smooth curves: if one adopts \(I_{\varnothing}=S\), the same degree convention gives \(\widehat\alpha(I_{\varnothing})=0\). This is an explicit convention for that boundary case, not an additional source hypothesis.
\subsection{Short English statement}
For the reduced set of singular points of a reduced plane curve, ask whether the asymptotic minimum degree per unit of imposed vanishing is always rational and whether it can be computed, particularly for Klein's arrangement of \(21\) lines. The intended ground field and the precise computational requirement remain unspecified.
\subsection{Sources}
\begin{itemize}
\item Supplied local UnsolvedMath registry record, permanent ID \(30004330\), title \emph{Waldschmidt Constants of Plane Curve Singular Sets}, source problem code \texttt{catalogue entry}, in the Oberwolfach Reports Open Problems collection. No individual proposer is supplied; the proposed year is recorded as \(2019\).
\item Supplied report citation: \emph{Mini-Workshop: Seshadri Constants} (2020), Oberwolfach Reports, pp.\ 3271--3313, DOI \texttt{10.4171/owr/2019/53}. \url{https://doi.org/10.4171/owr/2019/53}. The record describes these reports as documenting mathematical workshops and their open-problem discussions. The relationship between the citation year \(2020\) and the recorded proposed year \(2019\) is not explained.
\item Supplied literature source, labeled ``OWR report on Seshadri constants'': \url{https://ems.press/content/serial-article-files/46833}.
\item Supplied literature source, labeled ``Bauer et al., symmetric blowups and Waldschmidt constants'': \url{https://arxiv.org/abs/1609.08648}. No fuller bibliographic details are supplied.
\item Literature assessment dated exactly \texttt{2026-08-22}: the supplied curation labels the problem open and reports that rationality and computability in the Klein-arrangement test case remain unresolved in the primary literature, with later work providing estimates rather than a general solution. This is supplied dated curation derived from the OWR material, not an independently verified literature assessment. The record explicitly requests independent verification of time-sensitive status and cited claims.
\end{itemize}
\end{document}
